% Generated draft: requires mathematical and source review.
\documentclass[11pt]{article}
\usepackage[margin=25mm]{geometry}
\usepackage{fontspec}
\setmainfont{Latin Modern Roman}
\usepackage{amsmath,amssymb,mathtools,unicode-math}
\setmathfont{Latin Modern Math}
\usepackage{enumitem,needspace,xurl,hyperref}
\setlength{\parindent}{0pt}
\setlength{\parskip}{5pt}
\setlength{\emergencystretch}{3em}
\newcommand{\N}{\mathbb{N}}
\newcommand{\Z}{\mathbb{Z}}
\newcommand{\Q}{\mathbb{Q}}
\newcommand{\R}{\mathbb{R}}
\newcommand{\C}{\mathbb{C}}
\newcommand{\F}{\mathbb{F}}
\newcommand{\A}{\mathbb{A}}
\newcommand{\PP}{\mathbb{P}}
\newcommand{\E}{\mathbb{E}}
\newcommand{\eps}{\varepsilon}
\newcommand{\dd}{\,\mathrm{d}}
\begin{document}
\section*{UnsolvedMath ID: 30001940}
\textbf{Source problem code:} catalogue entry
\subsection{Definitions and mathematical statement}
The supplied record refers to a dissipative variational formulation of quantum fluid dynamics and to two energy functionals arising from it. Write \(E_1\) and \(E_2\) solely as labels for those functionals: their formulas, domains, and relation to the evolution are not supplied. The underlying space of measures, admissible transformations, and meaning of a dissipative Noether theory are likewise unspecified.

Subject to this unresolved source context, the problem consists of four research questions:
\begin{enumerate}
\item Are there other transformations, beyond those underlying the two referenced energy functionals, that yield new energy functionals for the dissipative variational formulation of quantum fluid dynamics?
\item How can a dissipative Noether theory be developed on the space of measures?
\item Can this theory be applied to classical fluid dynamics, and can it identify models possessing several energies?
\item Can this theory lead to an optimal-transport formulation of quantum mechanics?
\end{enumerate}
Here the unspecified transformations, the criterion for an energy functional to be new, the intended measure space, and the required content of each proposed theory or formulation must be recovered from the source context before these questions can be converted into precise mathematical assertions. No particular measure class, fluid model, boundary condition, regularity hypothesis, or quantifier over such choices is supplied, and none is imposed in this draft.
\subsection{Short English statement}
The source asks whether transformations can produce further quantum-fluid energy functionals, whether a dissipative Noether theory can be developed on measures and applied to classical fluids with several energies, and whether it can yield an optimal-transport formulation of quantum mechanics. The underlying functionals and mathematical framework are missing, so these remain research questions requiring source clarification.
\subsection{Sources}
\begin{itemize}
\item Local UnsolvedMath registry, supplied complete source record: permanent ID \(30001940\), title \emph{Dissipative Noether Theory and Quantum Fluid Energies}, Oberwolfach Reports Open Problems collection, exact source problem code \texttt{catalogue entry}. No individual proposer is identified; the proposed year is recorded as 2011.
\item Supplied source citation: \emph{Variational Methods for Evolution} (2012), Oberwolfach Reports, DOI 10.4171/owr/2011/55, pp.\ 3145--3216. Supplied source URL: \url{https://doi.org/10.4171/owr/2011/55}. The background attributes the questions to the report's workshop and open-problem discussions.
\item Supplied literature source, labeled ``original OWR report'': \url{https://doi.org/10.4171/owr/2011/55}.
\item Literature assessment checked on \textbf{2026-08-21}: the supplied curation labels the problem open and states that the report lists four questions concerning dissipative Noether theory, fluid dynamics, and optimal transport, with no complete resolution found. This is supplied dated curation, not an independently verified literature assessment; time-sensitive status and cited claims require independent verification.
\item The registry labels the statement ``exact'' and reports that its mathematical content was checked against the official Oberwolfach report, with notation, context, and question form normalized without changing the source-level claim. This verification description is supplied metadata and has not been independently confirmed here.
\end{itemize}
\end{document}
